\section{Introduction}

Quantitative real-time PCR (qPCR) is among the most widely used techniques
in molecular biology for quantifying nucleic acids and measuring gene
expression \cite{bustin2009miqe}. Current reporting guidance is provided by
MIQE~2.0 \cite{bustin2025miqe2}. A sound analysis involves several
interdependent steps: estimating primer amplification efficiency, choosing a
quantification strategy, testing groups statistically, and preparing figures
\cite{bustin2010miqe}. When target and reference efficiencies are comparable,
the $2^{-\Delta\Delta C_T}$ method \cite{livak2001analysis} is convenient;
when they differ, an efficiency-corrected or standard-curve-based approach is
required \cite{pfaffl2001new}. These choices, combined with the need for
rigorous statistics, make qPCR analysis demanding for many laboratory
researchers.

{\sloppy A range of tools cover different parts of this workflow
(Table~\ref{tab:tools}). Scripting-based implementations---our own
\textit{qPCRtools} \cite{li2022qpcrtools} and \textit{pcr}
\cite{ahmed2018pcr}---implement the calculation chain transparently but
require programming skill and a configured environment. Software distributed
with real-time instruments (for example QuantStudio Design \& Analysis and
CFX Maestro) offers a graphical workflow but is tied to that instrument.
Free stand-alone programs cover single steps: REST \cite{pfaffl2002rest}
computes efficiency-based relative expression for two-group comparisons with
a randomisation test, LinRegPCR
\cite{ramakers2003linregpcr,ruijter2009linregpcr} estimates baseline and
efficiency from raw fluorescence, and qBase \cite{hellemans2007qbase}
implements relative quantification with the geNorm stability measure
\cite{vandesompele2002accurate}. Free web applications---SATQPCR
\cite{rancurel2019satqpcr}, DAG Expression \cite{ballester2013dag} and
Auto-qPCR \cite{maussion2021autoqpcr}---integrate several steps but require
uploading the data to a remote server or deploying the web stack locally. The
gap is therefore not a missing algorithm but a missing delivery mode:
\emph{auditable numerical implementation} combined with a \emph{native,
cross-platform, zero-dependency} graphical workflow.
\par}

\begin{table}[h]
\centering
\footnotesize
\setlength{\tabcolsep}{2pt}
\caption{Delivery mode and functionality of representative qPCR analysis
tools. ``Free for non-commercial use'' means no purchase is required for
academic users; software distributed with an instrument is marked as not
free. ``Reference-gene stability'' denotes computational selection or ranking
of reference genes, not merely normalisation to a chosen reference gene.
$\bullet$~provided, $\circ$~not provided, $\circ$/$\bullet$~partial or
version-dependent (statistics were added in the later qBase+ release; vendor
software differs between instruments and versions). Compiled from the cited
papers, the vendors' documentation and the survey of Pabinger \textit{et al.}
\cite{pabinger2014survey}. \textsuperscript{a}~\textit{qPCRtools}
\cite{li2022qpcrtools}, \textit{pcr} \cite{ahmed2018pcr}, qpcR
\cite{ritz2008qpcr}, NormqPCR \cite{perkins2012normqpcr}, and the archived
Bioconductor package EasyqpcR \cite{lepape2017easyqpcr}.
\textsuperscript{b}~QuantStudio Design \& Analysis, CFX Maestro.
\textsuperscript{c}~SATQPCR \cite{rancurel2019satqpcr}, DAG Expression
\cite{ballester2013dag}, Auto-qPCR \cite{maussion2021autoqpcr}.}
\label{tab:tools}
\begin{tabular}{lcccccc}
\toprule
Tool & Free & No coding & Offline & Efficiency & Statistics & Ref.-gene stability \\
\midrule
R packages\textsuperscript{a} & $\bullet$ & $\circ$ & $\bullet$ & $\bullet$ & $\circ$/$\bullet$ & $\bullet$ \\
This application & $\bullet$ & $\bullet$ & $\bullet$ & $\bullet$ & $\bullet$ & $\bullet$ \\
REST & $\bullet$ & $\bullet$ & $\bullet$ & $\bullet$ & $\bullet$ & $\circ$ \\
LinRegPCR & $\bullet$ & $\bullet$ & $\bullet$ & $\bullet$ & $\circ$ & $\circ$ \\
qBase & $\bullet$ & $\bullet$ & $\bullet$ & $\bullet$ & $\circ$ & $\bullet$ \\
qBase+ & $\circ$ & $\bullet$ & $\bullet$ & $\bullet$ & $\bullet$ & $\bullet$ \\
Instrument software\textsuperscript{b} & $\circ$ & $\bullet$ & $\bullet$ & $\bullet$ & $\circ$/$\bullet$ & $\circ$ \\
Web apps\textsuperscript{c} & $\bullet$ & $\bullet$ & $\circ$ & $\bullet$ & $\bullet$ & $\bullet$ \\
\bottomrule
\end{tabular}
\end{table}

\subsection*{Statement of need}

The programming requirement is not a marginal obstacle. In our own teaching
and collaborative work, the users who most need a defensible qPCR analysis
are frequently those least able to install an R environment and write the
few lines of code the analysis requires: newly enrolled students, visiting
researchers, and wet-lab staff whose projects run for a single season. They
fall back either on manual spreadsheet calculations---difficult to audit and
easy to get wrong---or on whichever instrument software happens to be
installed. Our R package \cite{li2022qpcrtools} provides the algorithms, but
its use is bounded by the R environment itself.

As a pragmatic test of the workflow, we re-analysed the two real example
plates distributed with the 2022 \textit{qPCRtools} paper: one four-fold
dilution series and one separate treatment plate with four biological
replicates. Both analyses were completed through the graphical interface and
exported with a parameter record, without a spreadsheet or custom script
(Supplementary Table~S3). This is a reproducible published-data re-analysis,
not a controlled human-factors trial.

We therefore reimplemented the analytical core as a native desktop
application that implements the same three quantification methods while
removing every prerequisite except a double-click. \textit{qPCRtools} makes
three contributions: (i)~a C++/Qt architecture in which a tested core
algorithm library is decoupled from a web-based user interface; (ii)~a
reproducible validation against an independent base-R reference across all
quantification methods; and
(iii)~statistical and functional extensions beyond the R version
(non-parametric tests, effect sizes, outlier removal, quality-control
warnings, automatic reference-gene selection). Preparing the port also
revealed defects in result aggregation shared with the R package---the input
of the reference-gene stability measure, replicate handling in the efficiency
method, and the $\Delta\Delta C_t$ calibrator---as well as one specific to the
C++ code (multiple-comparison $P$-values); all were corrected and are covered
by regression tests.
